\section{Primitive tame data for the cubic ordinary Euler system}
\label{sec:cubic-primitive-euler}

The denominator permitted in AC-K29-ray is essential here.
The argument below chooses finitely many rational tame
correspondences and clears their denominators once.  It does not
claim an integral optimal-test-vector theorem or a bound for that
denominator as the modular form and its weight vary.
All automorphic representations in the local integrals below
are in their unitary normalization.  Consequently the classical
Rankin value at $s=e$ for weights $(2e+1,1)$ corresponds to the
local integral parameter $s_{\rm u}=0$.

\begin{lemma}[AC-K29-primitive-data: nonzero rational tame periods]
\label{AC-K29-primitive-data}
Let $f$ satisfy the hypotheses of AC-K29-ray, and put
$g=\theta_\chi$.  There is a finite set $S$ of rational primes,
disjoint from $\{\ell\}$, and algebraic tame test data
\[
 d_S=\bigotimes_{q\in S}(\Phi_q,W_{f,q},W_{g,q})
 \quad\text{in}\quad
 \bigotimes_{q\in S}
 \bigl(\mathcal S(\mathbb Q_q^2)\otimes\pi_{f,q}
                       \otimes\pi_{g,q}\bigr)
                                                        \tag{CEP.1}
\]
with the following properties.
The set $S$ contains every prime of ramification of $f$ or $g$.
At one prime $r\in S$ the Schwartz function satisfies
$\Phi_r(0)=0$ and $\int_{\mathbb Q_r^2}\Phi_r=0$.
For every $q\in S$ the normalized local period
\[
 B_q(d_q)=
 \left.
 \frac{Z(\Phi_q,W_{f,q},W_{g,q};s_{\rm u})}
      {L_q(\pi_f\times\pi_g,s_{\rm u})}
 \right|_{s_{\rm u}=0}
 \ne0.                                                   \tag{CEP.2}
\]
The data are finite linear combinations of tame translates
of the usual new vectors and torsion-coset Schwartz functions.
They may be defined over one finite extension of the coefficient
field.  Relative to any fixed lattices containing the new
vectors, a single power of the coefficient uniformizer clears
the denominators of all these data.
\end{lemma}
\begin{proof}
For each $q$, the local Rankin integrals generate
\[
 L_q(\pi_f\times\pi_g,s_{\rm u})
             \mathbb C[q^{s_{\rm u}},q^{-s_{\rm u}}].
                                                        \tag{CEP.3}
\]
This is Jacquet--Piatetski-Shapiro--Shalika (1983),
Theorem~2.7(ii), printed pp.~390--393, in the equal-rank
case.  In particular the normalized Laurent polynomials
cannot all vanish at $q^{s_{\rm u}}=1$.  Choose one datum
with nonzero value there.  Irreducibility of each local
generic representation implies that its new vector is
cyclic, so both chosen Whittaker vectors are finite sums
of translates of new vectors.

The choices can be algebraic.  The local representations of
the two algebraic modular forms have algebraic models; the
normalized new vectors and every fixed finite collection of
their translates have algebraic Whittaker values.  Integrating
over a compact coset contributes a rational Haar volume, and
the tails in the local integral are finite sums of
geometric series and their derivatives, with algebraic
ratios.  Thus their normalized Laurent
polynomials have algebraic coefficients.  Equivalently, the
algebraic test space spans the complex test space, so the
nonzero functional in (CEP.2) is already nonzero on an
algebraic pure tensor.  Only finitely many coefficients and
roots of unity are involved in a chosen datum; they lie in
one number field and hence in a finite coefficient extension.
This argument uses the rational local integral, not an
assertion about optimal integral vectors.

It remains to impose the global condition at the origin.
Choose a good prime $r\ne\ell$ at which
$\pi_{f,r}\not\simeq\pi_{g,r}^{\vee}$.
Such a prime exists outside any prescribed finite set.
Indeed, otherwise the characteristic polynomials of the
geometric Frobenius elements for
$V=M_f(e)$ and the finite-image representation of
$g^\vee$ would agree at all good primes.
Chebotarev and Brauer--Nesbitt would identify their
semisimplifications.  Their Hodge--Tate multisets are
$\{-e,e\}$ and $\{0,0\}$, respectively, which is
impossible for $e\ge1$.

The normalized functional $B_r$ is invariant under the
diagonal $\mathrm{GL}_2(\mathbb Q_r)$ at $s_{\rm u}=0$.
Put
\[
 \mathcal S_{00}=
 \{\Phi\in\mathcal S(\mathbb Q_r^2):
           \Phi(0)=0,\ \int\Phi=0\}.
\]
Evaluation and integration are independent linear
functionals on the Schwartz space.  With the row
action $(h\Phi)(x)=\Phi(xh)$, they identify its quotient
by $\mathcal S_{00}$ with
$1\oplus|\det|^{-1}$.
If $B_r$ vanished on
$\mathcal S_{00}\otimes\pi_{f,r}\otimes\pi_{g,r}$,
it would factor through this two-dimensional quotient.
The trivial-character contribution is
\[
 \operatorname{Hom}_{\mathrm{GL}_2(\mathbb Q_r)}
                 (\pi_{f,r}\otimes\pi_{g,r},1)=0.
\]
The equality follows from irreducibility and
$\pi_{f,r}\not\simeq\pi_{g,r}^{\vee}$.
The $|\det|^{-1}$ contribution is also zero:
the product central character of $\pi_{f,r}$ and
$\pi_{g,r}$ is $\eta_K^2=1$, whereas a central
scalar $zI$ acts on that extra character by $|z|^{-2}$.
This contradicts (CEP.3), so a datum in
$\mathcal S_{00}$ with $B_r(d_r)\ne0$ exists.
At the other primes use the choices already made, and use
the spherical data outside $S$.  The resulting global
Schwartz function and its Fourier transform both vanish
at the origin.  Fourier transform preserves
$\mathcal S_{00}$, since its value at zero is the
integral and its integral is the original value at zero.
Finally a finite collection of algebraic coefficients and
tame correspondences has a finite common denominator in
any fixed coefficient field.  This proves the last assertion.
\end{proof}

\begin{lemma}[AC-K29-primitive-reciprocity: fixed tame data at weight one]
\label{AC-K29-primitive-reciprocity}
Retain AC-K29-ray and fix finite algebraic tame data $d$
whose global Schwartz function vanishes at the origin.
Use these data in the universal ordinary Rankin--Eisenstein
construction, keeping the standard ordinary data at $\ell$.
Then the following statements hold.

First, after one fixed nonzero denominator, the construction
gives actual CM ray-class Euler systems $z_{\mathfrak n}(d)$
for the representations $U_a$ in (CER.2).  The denominator
is independent of the squarefree auxiliary ray modulus.
The norm relation is exactly (CER.3), after enlarging its
fixed exceptional set to include the support of $d$.
The classes satisfy the same away-$\ell$ local conditions
and the same $F^{--}=0$ condition as AC-K29-ray.

Second, for either ordinary CM refinement, the regulator of
the base class at $(k,k',j)=(2e-1,-1,e-1)$ is
\[
 \mathcal R_\ell(z_1(d))
  =u_{f,g,\ell,e}\,
       \frac{L^{\rm prim}(f,g,e)}{\Omega_{f,g,e}}\,
       \prod_{q\in S} B_q(d_q).
                                                        \tag{CEP.4}
\]
Here $\Omega_{f,g,e}$ is the nonzero classical period in
the Rankin interpolation formula, and
$u_{f,g,\ell,e}$ is nonzero.
The normalized algebraic right side is viewed in the
coefficient field using the fixed embeddings.
The regulator on the actual rank-one ordinary quotient is
the nonzero multiple (CER.5a) of the dual exponential.
In particular (CEP.4) is an assertion about the actual
weight-one specialization, even when the usual
imprimitive value at that point is zero.
\end{lemma}
\begin{proof}
\proofstep{1. Finite tame operations on the actual universal class.}
Use the canonical equivariant Siegel-unit map
\[
 \mathcal S_0(\mathbb A_f^2,\mathbb Q)
       \longrightarrow H^1_{\rm mot}(Y_{\mathrm{GL}_2},\mathbb Q(1)),
\]
followed by the diagonal pushforward.  In the explicit
infinite-level normalization this is Groutides,
arXiv:2407.01377v3, Theorem~5.2.3,
Proposition~5.2.2 and Definition~5.2.4, printed pp.~18--19.
The volume in Proposition~5.2.2 is part of the
pushforward, so that this is an equivariant map rather
than an unspecified choice of finite-level normalization.
Use its $c$-smoothed integral étale realization, as in
the same paper, \S5.2.5, and KLZ
arXiv:1503.02888v6, Theorem~6.3.4.

Here is a direct finite-level description of all the data
being used.  A locally constant compactly supported
Schwartz function which vanishes at zero is a finite sum
of characteristic functions of nonzero rational lattice
cosets.  After clearing denominators a coset has the form
$(a,b)+N\widehat{\mathbb Z}^{\,2}$, with its zero coset
absent.  Put $d_0=\gcd(a,b,N)$ and $M=N/d_0$.
After division by $d_0$, its row is primitive modulo $M$.
An element of $\mathrm{SL}_2(\mathbb Z/M\mathbb Z)$ takes
that row to $(0,1)$.  Thus the coset function is a tame
translate and dilation of
$\mathbf1_{(0,1)+M\widehat{\mathbb Z}^{\,2}}$.
For data away from $\ell$, all these operations can be
taken in $\mathrm{GL}_2(\mathbb A_f^{(\ell)})$ and all
denominators and levels are prime to $\ell$.
The equivariance transfers the diagonal translate to the
two modular test vectors.  A vector in either local
Whittaker model is a finite sum of translates of its
new vector, so the entire construction is a finite linear
combination of standard torsion-section constructions,
finite pullbacks and finite pushforwards away from $\ell$.

Apply exactly the fixed first-factor moment and ordinary
quotient, and the free CM quotient, of AC-K29-ray.
In particular the second factor at weight one is the
actual ordinary Hida specialization.  No negative-degree
algebraic moment is used.  Each newly introduced
correspondence is at a fixed tame prime and fixed level;
its degree and its coefficient denominator are independent
of the ray modulus and of the cyclotomic layer.
Their finite sum therefore has one common denominator.
In detail, let $\pi_{\mathfrak n}$ denote the actual
integral CM quotient map at the varying ray modulus
used in AC-K29-ray.  Choose once a smaller compact
open subgroup at the primes in $S$ on which all our
vectors and Schwartz functions are fixed.
Every chosen second-factor operator is a fixed
pullback followed by a fixed pushforward from this
enhanced tame level to the level on which
$\pi_{\mathfrak n}$ is defined.  The new coefficient
map is $\pi_{\mathfrak n}\circ T$, not an assumed
scalar action of $T$ on the rank-two quotient.
The compact open subgroups outside $S$ retain
their original dependence on $\mathfrak n$.
Thus the indices and coefficients in $T$ are fixed,
whereas $\pi_{\mathfrak n}$ has the same integral
bound as in the existing CM construction.

\proofstep{2. The norm and local-condition assertions.}
At a new auxiliary prime $\mathfrak q$, none of the
fixed tame operators has support at $\mathfrak q$.
They commute with the pushforward and the good-prime
Hecke operators in the proof of (CER.3).
Thus that proof applies to each summand and
then to their sum with exactly the same polynomial
$P_{\mathfrak q}([\mathfrak q])$.
More precisely, before taking the quotient the
norm correspondence at $\mathfrak q$ commutes
with $T$, since their finite adelic supports are
disjoint.  The original compatibility of
$\pi_{\mathfrak n}$ with this norm correspondence
therefore implies the compatibility of
$\pi_{\mathfrak n}\circ T$ with the same
coefficient norm map.  This proves the assertion
on the actual varying CM ray modules; it requires
neither a new multiplicity-one quotient at the
enhanced level nor a denominator depending on
$\mathfrak n$.
This is also the mechanism in the full proof of
Groutides Theorem~5.2.7 and Corollary~5.2.10:
the data outside the new prime are fixed and arbitrary,
and the local norm identity is transported by the
equivariant Rankin--Eisenstein map.
The separate CM ray refinement is the one in
AC-K29-ray; no replacement by a merely cyclotomic
Euler system is made.

The operators commute with the ordinary projectors at
$\ell$ and with the coefficient filtrations.
The proof of KLZ Proposition~8.1.7, including its
vanishing of the $F^{--}$ projection, is therefore
unchanged.  Before specialization these are still
cyclotomic universal norms.  The same universal-norm
argument used in AC-K29-ray gives the Bloch--Kato
condition away from $\ell$.

\proofstep{3. The scalar reciprocity identity with the finite tame operators.}
We explain the extension of the scalar reciprocity
calculation, because simply cancelling a vanishing
Euler factor would not identify the endpoint value.
Reduce first to a standard torsion-coset function by
Step~1.  The second modular form is then a finite tame
translate of its usual vector at a fixed larger level,
and the first functional is its usual eigenform
functional composed with a finite tame correspondence.

The required classical regulator calculation is explicitly
valid for these vectors.  In KLZ,
arXiv:1501.03289v2, Remark~6.4.5, the second form may
be any holomorphic form at the chosen level with the
same Hecke eigenvalues away from that level.
The full calculation in Propositions~6.3.1--6.3.3,
Corollary~6.4.7, Propositions~6.5.4--6.5.6,
Corollary~6.5.7 and Theorem~6.5.9 is linear in that
form, the first de Rham functional, and the Eisenstein
symbol.  More explicitly, its initial step is the
projection formula
\[
 \langle\Delta_*\mathrm{Eis},\,\eta\boxtimes\omega\rangle
       =\langle\mathrm{Eis},\,\eta\cup\omega\rangle.
                                                        \tag{CEP.5}
\]
The first class is arbitrary until its unit-root
eigenvalue is selected.  A tame correspondence commutes
with crystalline Frobenius, and hence preserves that
selection.  The subsequent operations are cup products,
the fixed Clebsch--Gordan map, the unit-root splitting,
and the ordinary projector.  They commute with the
finite pullbacks and pushforwards just described.
The Eisenstein series is the one attached to the
chosen torsion section, with the same normalization
on both sides of (CEP.5).
Thus the displayed calculation proves the same scalar
reciprocity formula for every summand of Step~1.

Ohta's comparison and the first-factor eigenfunctional
specialize as in KLZ Proposition~10.1.1; the correction
in \S13 keeps the functional regular and nonzero at
the fixed positive-weight $f$, even if its residual
representation is reducible.  The actual CM rank-one
functional is the canonical isomorphism in
Proposition~10.1.1(1).
The torsion-free scalar-function argument in the full
proof of Theorem~10.2.2 consequently applies to the
finite tame operators as well.
This is a comparison of scalar regulator functions,
not a claim that a dense set of specializations detects
an arbitrary Galois cohomology class.
Throughout this step the classes and their coefficient
maps are the ordinary KLZ ones.
The scalar analytic identities from Loeffler used
next do not identify an ordinary Hida cohomology
class with a Coleman cohomology class.

\proofstep{4. Critical interpolation at the actual CM weight-one point.}
For completeness, the scalar interpolation at weight
one follows from the explicit two-slice argument of
Loeffler, arXiv:1704.04049v2, \S\S4--6.
It continues to apply to the vectors just described:
the second family in \S4 is expressly allowed to be
neither normalized nor cuspidal, and the construction
is linear in the first functional.  A fixed tame
operator is a bounded operator on these ordinary or
overconvergent families at a fixed larger tame level.
It also commutes with $\ell$-depletion and with the
Gauss--Manin operations, with the usual determinant
power included in the slash action.

Here is the algebraic identity responsible for the
continuation.  For each reduced standard torsion section,
write $A(d,d')$ for its fixed periodic Fourier coefficient
(namely $\zeta_M^d+(-1)^{\kappa_1-\kappa_2}\zeta_M^{-d}$
for the reduced standard torsion section).
It is constant in $d'$ on a fixed parity component.
The two
Eisenstein families have nonconstant coefficients
\[
 E_\kappa:\ \sum_{dd'=n,\ \ell\nmid n}
                  d^{\kappa-1}A(d,d'),\qquad
 F_\kappa:\ \sum_{dd'=n,\ \ell\nmid n}
                  (d')^{\kappa-1}A(d,d').
                                                        \tag{CEP.6}
\]
For the standard torsion point these are precisely
Loeffler's Lemma~3.2.  The diagonal tame translates
have already been transferred to the modular test
vectors, so no additional assertion about translating
the differential operator in (CEP.6) is needed.
Finite torsion combinations give the same identity
by linearity.
If
$\tau+\tau'=\kappa_1-\kappa_2-1$, multiplication by
$n^\tau$ and $n^{\tau'}$ gives
\[
 \theta^\tau E_{\kappa_1-\kappa_2-2\tau}
       =\theta^{\tau'}F_{\kappa_1-\kappa_2-2\tau'}.
                                                        \tag{CEP.7}
\]
Multiplying by the same second family and applying the
same first functional preserves this equality.
The tame-level powers in the two normalizations agree
on this intersection as well.

For an integer $\tau=t\ge0$, the geometric-specialization
calculation of Step~3 identifies the scalar regulator
with the first slice on a dense set of positive-weight
arithmetic points.  Indeed, writing
$\kappa_1=k+2,\ \kappa_2=k'+2$ and
$\sigma=\kappa_1-1-t$, its geometric exponent is
$j=k-t$; the inequalities
$k\ge t,\ k'\ge k-t$ supply a dense set in the two
weight variables.  One can also require $k+2>2t$
to avoid every projector pole.  The identity theorem
gives equality on that slice.
Equation (CEP.7) then gives equality on its
intersections with the second slices.  Their union is
dense by Loeffler Lemma~5.3, so the argument of
Corollary~6.2 gives equality on each second slice.
The CM ordinary family is available on the full
weight component, exactly as required there.

At our point use the second slice
$\tau'=\sigma-\kappa_2=e-1$.
The near-overconvergence degree is $e-1$ and
\[
 \kappa_1=2e+1>2(e-1),\qquad
 1=\kappa_2\le\kappa_1-1-(e-1)=e+1.
                                                        \tag{CEP.8}
\]
Thus Proposition~4.1 identifies the specialization
with the classical holomorphic Rankin period, with
no pole in the projector.  This proves the endpoint
identity directly, including when the imprimitive
period is zero.

\proofstep{5. Unfolding and the nonzero local scalar.}
We first fix the Fourier convention in this step.
This can be checked on the standard torsion section,
not merely from the isomorphism type of a principal
series.  For
$\Phi_M=\mathbf1_{(0,b)+M\widehat{\mathbb Z}^{\,2}}$,
its symplectic Fourier transform is supported on
$M^{-1}\widehat{\mathbb Z}^{\,2}$ and, on the
rational lattice in that support, is
\[
 \widehat\Phi_M(m/M,n/M)
          =M^{-2}\exp(2\pi i m b/M).
\]
Here the finite additive character is chosen opposite
to the usual real one; reversing that convention
replaces $b$ by $-b$ on both sides.
Substituting this formula into the Godement lattice
sum, and making the fixed dilation by $M$, gives
the series
$F^{\rm an}_{t,s,b}$ of KLZ2015
Definition~4.4.3, with its displayed nonzero volume
and determinant powers.
Proposition~4.4.5 identifies this series with the
actual absolute-Hodge Eisenstein regulator; in
degree zero it gives
$F^{\rm an}_{0,1,b}=-2\log|g_{0,b/M}|$.
Definition~4.3.1 and Proposition~4.3.3 give its
de Rham derivative, the same Eisenstein series
used in the holomorphic period.
Finite coset reduction and equivariance therefore
give the identification for every datum in Step~1.
In the Godement notation this is precisely
$f_{\widehat\Phi}(1,1,1)$, as in Loeffler,
arXiv:2011.15106, Proposition~3.3(ii).
It is therefore not legitimate to replace
$\widehat\Phi$ by $\Phi$ in an unnormalized integral.
Write $\widetilde W_f,\widetilde W_g$ for the
contragredient Whittaker vectors obtained by the
usual long-Weyl transpose-inverse operation.
The exact normalized local functional equation,
JPSS Theorem~2.7(iii), is
\[
 \frac{Z(\widehat\Phi,\widetilde W_f,\widetilde W_g;1-s)}
      {L_q(\pi_f^\vee\times\pi_g^\vee,1-s)}
  =\varepsilon_q^\sharp(s)
       \frac{Z(\Phi,W_f,W_g;s)}
            {L_q(\pi_f\times\pi_g,s)} .
                                                        \tag{CEP.8a}
\]
Here $\varepsilon_q^\sharp$ includes the
central-character sign in that theorem and is a
nonzero monomial in $q^{-s}$.
The contragredient convention for the first
de Rham functional is compatible with this formula:
$\pi_f^\vee=\pi_f\otimes\eta_K$ and
$\pi_g\otimes\eta_K\simeq\pi_g$.
An $\eta_K$ factor can therefore be transferred
between the two Whittaker vectors under the
diagonal integral without changing the Rankin
representation.  These are fixed invertible
identifications of test spaces, not scalar
specialization maps which could vanish.
A symplectic instead of ordinary Fourier transform
only inserts a fixed long-Weyl translate in this
formula.  All these transforms and translates are
the fixed tame operations already allowed above.
Both sides of (CEP.8a), after normalization, are
regular at $s=0$; its right side there is
$\varepsilon_q^\sharp(0)B_q(d_q)$.
This is how the $s=1$ Godement convention is
converted to the $s_{\rm u}=0$ convention of
(CEP.2).  In particular no unnormalized
intertwining scalar is assumed nonzero at a
reducible principal-series endpoint.
The auxiliary datum of the first lemma lies in
$\mathcal S_{00}$, so either Fourier convention
satisfies the required global origin condition.

Unfold the classical period with the chosen global
Schwartz function and modular vectors.  At every
finite place use (CEP.8a) if the Godement
Fourier convention is used.  Factoring out the primitive local
$L$-factors therefore gives
$L^{\rm prim}(f,g,e)\prod_{q\in S}B_q(d_q)$,
up to the nonzero epsilon and normalization factors.
If the unfolding is written at the complementary
unitary parameter $1$, the global functional
equation makes the same conversion.  Its gamma
factors at $0$ and $1$ are finite and nonzero,
as in AC-K29-bad, so this introduces another
nonzero factor in the algebraic Rankin normalization.
This is the same Rankin unfolding as in Loeffler
Proposition~2.10 and Appendix~A; changing the fixed
tame sections changes precisely these local
integrals.  The archimedean vector and the data at
$\ell$ were kept fixed.  Their factors are the
factorials and ordinary Euler factors of AC-K29-lower,
which are nonzero at (CEP.8).

Choose the smoothing integer $c$ congruent to one
modulo the common tame level and with
$c^2\not\equiv1\pmod\ell$.  The Chinese remainder
theorem permits this.  The smoothing operator on
our datum then has scalar $c^2-1$ at the point in
question.  It is nonzero.  All other normalization
factors are nonzero tame volumes, determinant powers,
the nonzero first eigenfunctional, and the canonical
CM functional.  This proves (CEP.4).

Finally the local big logarithm is functorial on
the unramified ordinary quotient.  Specializing the
actual CM quotient at weight one and twisting by
$j=e-1$ gives exactly (CER.5a), by the full construction
in KLZ Theorems~8.2.3 and~8.2.8.
The tame operators have changed neither this quotient
nor its Frobenius eigenvalues.  Its scalar is nonzero
by AC-K29-lower's local calculation.  This completes
the proof.
\end{proof}

\begin{theorem}[AC-K29-primitive: removal of the shared bad-factor condition]
\label{AC-K29-primitive}
Assume all the hypotheses of AC-K29-lower except (CER.5).
Then
\[
       H^1_f(K,V\chi)=H^1_f(K,V\chi^2)=0.
                                                        \tag{CEP.9}
\]
In particular the matching local type of AC-K29-local-types
at $269$ does not require an additional imprimitive
nonvanishing assumption for this coefficient-field theorem.
\end{theorem}
\begin{proof}
Choose the tame data of AC-K29-primitive-data and
apply AC-K29-primitive-reciprocity to both ordinary CM
refinements.  AC-K29-bad proves that the primitive
complex value is nonzero.  Every factor in (CEP.4)
is therefore nonzero, so both actual base classes
have nonzero singular localization at $\ell$.
They still have zero $F^{--}$ projection and lie in
the two distinct CM Frobenius lines.  Hence the two
singular images span the singular local space.

The Poitou--Tate argument (CER.7)--(CER.8) identifies
the Bloch--Kato group with the group strict at both
primes above $\ell$.  The ray Euler systems have a
single fixed denominator, exactly the required norm
polynomials, and the away-$\ell$ local conditions.
One of their base classes is nonzero.
The rank-two $K$-Euler-system argument used in the
last paragraph of AC-K29-lower therefore kills this
strict group over the coefficient field, using the
retained derived-image hypothesis of AC-K29-split.
This proves (CEP.9) for both conjugate characters.
\end{proof}

\paragraph{Remaining scope and geometric dictionary.}
The proof uses modular-curve Siegel units, their diagonal
pushforward, ordinary Iwasawa coefficient sheaves, and the
actual CM ray-class quotient.  These are the same étale
realizations and norm maps as in AC-K29-ray.
The choice of additional tame data is finite for each
fixed modular form, weight and coefficient prime.
It supplies no uniform estimate for an elliptic-unit
specialization index, no conclusion at coefficient
primes outside AC-K29-ray's ordinary split
distinguished scope, and no assertion that every point
of a new crystalline component is ordinary.
The rank-four rational Euler-system image obstruction
AC-K30 is not bypassed by an unproved image assertion:
the strict group is controlled by the rank-two system
over $K$.

\paragraph{Source audit.}
JPSS Theorem~2.7(i)--(iii) is used in its equal-rank
case; its complete printed arguments on pp.~391--393
and \S\S2.9--2.10, pp.~395--398, were read,
including the common denominator, ideal property,
generator normalization, Fourier covariance and
the nonzero monomial epsilon factor.
The complete applied proofs of Groutides
Theorem~5.2.7 and Corollary~5.2.10 were read;
their upstream representation-theoretic norm theorem
is retained as a source input.
Loeffler--Zerbes, arXiv:2312.04665v2,
Proposition~3.4.1 and Remark~3.4.2 provide the analogous
fixed-denominator primitive Asai construction.
They are used as a construction check, not as a
Rankin theorem or as a weight-one coefficient comparison.
The exact Rankin and weight-one extensions used here
are proved in the second lemma from the applied
KLZ regulator calculation and the displayed
two-slice identities.
